% !TeX root = ../shcmthesis-example.tex

\part{\LaTeX{} 使用}

\chapter{\LaTeX{} 介绍}

\textbf{声明：以下文字节选复制自\url{https://oi-wiki.org/tools/latex/}，并根据本文情况稍作修改}

\section{什么是 \LaTeX{}}

LaTeX（读作/ˈlɑːtɛx/或/ˈleɪtɛx/）是一个让你的文档看起来更专业的排版系统，而不是文字处理器。它尤其适合处理篇幅较长、结构严谨的文档，并且十分擅长处理公式表达。它是免费的软件，对大多数操作系统都适用。

LaTeX 基于 TeX（Donald Knuth 在 1978 年为数字化排版设计的排版系统）。TeX 是一种电脑能够处理的低级语言，但大多数人发现它很难使用。LaTeX 正是为了让它变得更加易用而设计的。目前 LaTeX 的版本是 LaTeX 2e。

如果你习惯于使用微软的 Office Word 处理文档，那么你会觉得 LaTeX 的工作方式让你很不习惯。Word 是典型的「所见即所得」的编辑器，你可以在编排文档的时侯查看到最终的排版效果。但使用 LaTeX 时你并不能方便地查看最终效果，这使得你专注于内容而不是外观的调整。

一个 LaTeX 文档是一个以 .tex 结尾的文本文件，可以使用任意的文本编辑器编辑，比如 Notepad，但对于大多数人而言，使用一个合适的 LaTeX 编辑器会使得编辑的过程容易很多。在编辑的过程中你可以标记文档的结构。完成后你可以进行编译——这意味着将它转化为另一种格式的文档。它支持多种格式，但最常用的是 PDF 文档格式。

\section{章节}

文档可以分为篇（Part）、章（Chatpers）、节（Sections）小节（Subsections）和小小节（Subsubsections），使用命令如下：

\begin{python}
\part{...}
\chapter{...}
\section{...}
\subsection{...}
\subsubsection{...}
\end{python}

\section{文字处理}

LaTeX 有多种不同的字体效果，在此列举一部分：

\begin{python}
\textit{words in italics} 
\textsl{words slanted} 
\textsc{words in smallcaps} 
\textbf{words in bold} 
\texttt{words in teletype} 
\textsf{sans serif words} 
\textrm{romanwords} 
\underline{underlined words}
\end{python}

效果如下：

\textit{words in italics} 

\textsl{words slanted} 

\textsc{words in smallcaps} 

\textbf{words in bold} 

\texttt{words in teletype} 

\textsf{sans serif words} 

\textrm{romanwords} 

\underline{underlined words}

其中，这些命令都可以通过快捷键打出，选中一句文字后按下快捷键即可，例如\textbackslash textit\{\} 为Ctrl+i，\textbackslash textbf\{\} 为Ctrl+b，
如果使用苹果电脑则将Ctrl替换为Command。

\section{段落}

\subsection{段落换行}
在LaTex中段落需要间隔一行以进行换行，例如：

\begin{python}
段落一
段落二

段落三
\end{python}

显示效果为：

段落一
段落二

段落三

\subsection{段落缩进}

LaTeX 默认每个章节第一段首行顶格，之后的段落首行缩进。如果想要段落顶格，在要顶格的段落前加 \noindent 命令即可。

\section{列表}

LaTeX 支持两种类型的列表：有序列表（enumerate）和无序列表（itemize）。列表中的元素定义为\textbackslash item。列表可以有子列表。

\rightarrow 输入下面的内容来生成一个有序列表套无序列表：

\begin{python}
\begin{enumerate}
	\item First thing
	
	\item Second thing
	\begin{itemize}
		\item A sub-thing
		
		\item Another sub-thing
	\end{itemize}
	
	\item Third thing
\end{enumerate}
\end{python}

\rightarrow 编译并核对 PDF 文档。

列表长这样：

\begin{enumerate}
	\item First thing
	
	\item Second thing
	\begin{itemize}
		\item A sub-thing
		
		\item Another sub-thing
	\end{itemize}
	
	\item Third thing
\end{enumerate}

可以使用方括号参数来修改无序列表头的标志。例如，\textbackslash item[-] 会使用一个杠作为标志，你甚至可以使用一个单词，比如 \textbackslash item[One]。

下面的代码：

\begin{python}
\begin{itemize}
	\item[-] First thing
	
	\item[+] Second thing
	\begin{itemize}
		\item[Fish] A sub-thing
		
		\item[Plants] Another sub-thing
	\end{itemize}
	
	\item[Q] Third thing
\end{itemize}
\end{python}


生成的效果为：

\begin{itemize}
	\item[-] First thing
	
	\item[+] Second thing
	\begin{itemize}
		\item[Fish] A sub-thing
		
		\item[Plants] Another sub-thing
	\end{itemize}
	
	\item[Q] Third thing
\end{itemize}

\section{注释和空格}

我们使用 \% 创建一个单行注释，在这个字符之后的该行上的内容都会被忽略，直到下一行开始。快捷键为Ctrl+t，若使用苹果电脑则为Command+t。

下面的代码：

\begin{python}
It is a truth universally acknowledged % Note comic irony
in the very first sentence , 
that a single man in possession of a good fortune,
must be in want of a wife.
\end{python}

生成的结果为：

It is a truth universally acknowledged % Note comic irony
in the very first sentence , that a single man in possession of a good fortune,
must be in want of a wife.

多个连续空格在 LaTeX 中被视为一个空格。多个连续空行被视为一个空行。空行的主要功能是开始一个新的段落。通常来说，LaTeX 忽略空行和其他空白字符，两个反斜杠（\textbackslash\textbackslash）可以被用来换行。

如果你想要在你的文档中添加空格，你可以通过添加~从而添加空格。
或者 \textbackslash vaspace\{...\} 的命令，这样可以添加竖着的空格，高度可以指定。如 \textbackslash vspace\{12pt\} 会产生一个空格，高度等于 12pt 的文字的高度。

\section{脚注}

在需要添加脚注的地方插入\textbackslash footnote{脚注内容}，则会自动在页脚部分插入脚注内容。

例如：

\begin{python}
文字内容一\footnote{注释一}和文字内容二\footnote{注释二}
\end{python}

效果如下：

文字内容一\footnote{注释一}和文字内容二\footnote{注释二}

可以看到在本页底部已经自动按编号添加了注释内容。

\section{特殊字符}

下列字符在 LaTeX 中属于特殊字符：

\begin{python}
#
$
%
^
&
_
{
}
~
\
\end{python}

为了使用这些字符，我们需要在他们前面添加反斜杠进行转义：

\begin{python}
\#
\$
\%
\^{}
\&
\_
\{
\}
\~{}
\end{python}

\section{表格}

表格（tabular）命令用于排版表格，表应具有自明性。

\subsection{普通表格}

LaTeX 默认表格是没有横向和竖向的分割线的——如果你需要，你得手动设定。LaTeX 会根据内容自动设置表格的宽度。

下面的代码可以创一个表格：

\begin{python}
\begin{tabular}{...}
\end{tabular}
\end{python}
	
例如：

\begin{python}
\begin{center}
\begin{tabular}{|r|l|}
	\centering
	\hline
	8              & here's \\
	\cline 86 & stuff  \\
	\hline
	\hline
	2008           & now    \\
	\hline
\end{tabular}
\end{center}
\end{python}

效果如下：

\begin{center}
\begin{tabular}{|r|l|}
	\hline
	8              & here's \\
	\cline{2-2}
	86 & stuff  \\
	\hline
	2008           & now    \\
	\hline
\end{tabular}
\end{center}

l 表示一个左对齐的列；

r 表示一个右对齐的列；

c 表示一个向中对齐的列；

| 表示一个列的竖线；

例如，\{lll\} 会生成一个三列的表格，并且保存向左对齐，没有显式的竖线；\{|l|l|r|\} 会生成一个三列表格，前两列左对齐，最后一列右对齐，并且相邻两列之间有显式的竖线。

表格的数据在 \textbackslash begin\{tabular\} 后输入：

\& 用于分割列；

\textbackslash \textbackslash 用于换行；

\textbackslash hline 表示插入一个贯穿所有列的横着的分割线；

\textbackslash cline\{1-2\} 会在第一列和第二列插入一个横着的分割线。

最后使用 \textbackslash end\{tabular\} 结束表格。

\subsection{三线表}

为使表格简洁易读，可以采用三线表，如表~\ref{tab:three-line}。
三条线可以使用 \pkg{booktabs} 宏包提供的命令生成。

例如：

\begin{python}
\begin{table}
	\centering
	\caption{三线表示例}
	\begin{tabular}{ll}
		\toprule
		文件名          & 描述                         \\
		\midrule
		shcmthesis.dtx   & 模板的源文件，包括文档和注释 \\
		shcmthesis.cls   & 模板文件                     \\
		shcmthesis-*.bst & BibTeX 参考文献表样式文件    \\
		\bottomrule
	\end{tabular}
	\label{tab:three-line}
\end{table}
\end{python}

效果如下：

\begin{table}
	\centering
	\caption{三线表示例}
	\begin{tabular}{ll}
		\toprule
		文件名          & 描述                         \\
		\midrule
		shcmthesis.dtx   & 模板的源文件，包括文档和注释 \\
		shcmthesis.cls   & 模板文件                     \\
		shcmthesis-*.bst & BibTeX 参考文献表样式文件    \\
		\bottomrule
	\end{tabular}
	\label{tab:three-line}
\end{table}

表格如果有附注，尤其是需要在表格中进行标注时，可以使用 \pkg{threeparttable} 宏包。
研究生要求使用英文小写字母 a、b、c……顺序编号，本科生使用圈码 ①、②、③……编号。

例如：

\begin{python}
\begin{table}
	\centering
	\begin{threeparttable}[c]
		\caption{带附注的表格示例}
		\label{tab:three-part-table}
		\begin{tabular}{ll}
			\toprule
			文件名                 & 描述                         \\
			\midrule
			shcmthesis.dtx\tnote{a} & 模板的源文件，包括文档和注释 \\
			shcmthesis.cls\tnote{b} & 模板文件                     \\
			shcmthesis-*.bst        & BibTeX 参考文献表样式文件    \\
			\bottomrule
		\end{tabular}
		\begin{tablenotes}
			\item [a] 可以通过 xelatex 编译生成模板的使用说明文档；
			使用 xetex 编译 \file{shcmthesis.ins} 时则会从 \file{.dtx} 中去除掉文档和注释，得到精简的 \file{.cls} 文件。
			\item [b] 更新模板时，一定要记得编译生成 \file{.cls} 文件，否则编译论文时载入的依然是旧版的模板。
		\end{tablenotes}
	\end{threeparttable}
\end{table}
\end{python}

效果如下：

\begin{table}
	\centering
	\begin{threeparttable}[c]
		\caption{带附注的表格示例}
		\label{tab:three-part-table}
		\begin{tabular}{ll}
			\toprule
			文件名                 & 描述                         \\
			\midrule
			shcmthesis.dtx\tnote{a} & 模板的源文件，包括文档和注释 \\
			shcmthesis.cls\tnote{b} & 模板文件                     \\
			shcmthesis-*.bst        & BibTeX 参考文献表样式文件    \\
			\bottomrule
		\end{tabular}
		\begin{tablenotes}
			\item [a] 可以通过 xelatex 编译生成模板的使用说明文档；
			使用 xetex 编译 \file{shcmthesis.ins} 时则会从 \file{.dtx} 中去除掉文档和注释，得到精简的 \file{.cls} 文件。
			\item [b] 更新模板时，一定要记得编译生成 \file{.cls} 文件，否则编译论文时载入的依然是旧版的模板。
		\end{tablenotes}
	\end{threeparttable}
\end{table}

如某个表需要转页接排，可以使用 \pkg{longtable} 宏包，需要在随后的各页上重复表的编号。
编号后跟表题（可省略）和“（续）”，置于表上方。续表均应重复表头。

例如：

\begin{python}
\begin{longtable}{cccc}
	\caption{跨页长表格的表题}
	\label{tab:longtable} \\
	\toprule
	表头 1 & 表头 2 & 表头 3 & 表头 4 \\
	\midrule
	\endfirsthead
	\caption*{续表~\thetable\quad 跨页长表格的表题} \\
	\toprule
	表头 1 & 表头 2 & 表头 3 & 表头 4 \\
	\midrule
	\endhead
	\bottomrule
	\endfoot
	Row 1  & & & \\
	Row 2  & & & \\
	Row 3  & & & \\
	Row 4  & & & \\
	Row 5  & & & \\
	Row 6  & & & \\
	Row 7  & & & \\
	Row 8  & & & \\
	Row 9  & & & \\
	Row 10 & & & \\
\end{longtable}
\end{python}

效果如下：

\begin{longtable}{cccc}
	\caption{跨页长表格的表题}
	\label{tab:longtable} \\
	\toprule
	表头 1 & 表头 2 & 表头 3 & 表头 4 \\
	\midrule
	\endfirsthead
	\caption*{续表~\thetable\quad 跨页长表格的表题} \\
	\toprule
	表头 1 & 表头 2 & 表头 3 & 表头 4 \\
	\midrule
	\endhead
	\bottomrule
	\endfoot
	Row 1  & & & \\
	Row 2  & & & \\
	Row 3  & & & \\
	Row 4  & & & \\
	Row 5  & & & \\
	Row 6  & & & \\
	Row 7  & & & \\
	Row 8  & & & \\
	Row 9  & & & \\
	Row 10 & & & \\
\end{longtable}

\subsection{表格引用}

通过在表格中加入\textbackslash label\{tab\_name\}来定义表格名称，在文本中使用\textbackslash ref\{tab\_name\}来引用。

建议在\textbackslash ref前添加空格~以示美观。

例如：

\begin{python}
此处引用表~\ref{tab:longtable}。
\end{python}

效果为：

此处引用表~\ref{tab:longtable}。

\section{图片}

本章介绍如何在 LaTeX 文档中插入图片。

图片应当是 PDF，PNG，JPEG 或者 GIF 文件。
建议矢量图片使用 PDF 格式，比如数据可视化的绘图；
照片应使用 JPG 格式；
其他的栅格图应使用无损的 PNG 格式。
注意，LaTeX 不支持 TIFF 格式；EPS 格式已经过时。

\subsection{单张图片}

图片通常在 \env{figure} 环境中使用 \cs{includegraphics} 插入，如图~\ref{fig:example} 的源代码。

例如：

\begin{python}
\begin{figure}
	\centering
	\includegraphics[width=0.5\linewidth]{example-image-a.pdf}
	\caption*{国外的期刊习惯将图表的标题和说明文字写成一段，需要改写为标题只含图表的名称，其他说明文字以注释方式写在图表下方，或者写在正文中。}
	\caption{示例图片标题}
	\label{fig:example}
\end{figure}
\end{python}

效果如下：

\begin{figure}
	\centering
	\includegraphics[width=0.5\linewidth]{example-image-a.pdf}
	\caption*{国外的期刊习惯将图表的标题和说明文字写成一段，需要改写为标题只含图表的名称，其他说明文字以注释方式写在图表下方，或者写在正文中。}
	\caption{示例图片标题}
	\label{fig:example}
\end{figure}

[h] 是位置参数，h 表示把图表近似地放置在这里（如果能放得下）。有其他的选项：t 表示放在在页面顶端；b 表示放在在页面的底端；p 表示另起一页放置图表。你也可以添加一个 ! 参数来强制放在参数指定的位置（尽管这样排版的效果可能不太好）。

\textbackslash centering 将图片放置在页面的中央。如果没有该命令会默认左对齐。使用它的效果是很好的，因为图表的标题也是居中对齐的。

\textbackslash includegraphics\{...\} 命令可以自动将图放置到你的文档中，图片文件应当与 TeX 文件放在同一目录下。

width=1\textbackslash textwidth 是一个可选的参数，它指定图片的宽度——与文本的宽度相同。宽度也可以以厘米为单位。你也可以使用 [scale=0.5] 将图片按比例缩小（示例相当于缩小一半）。

\textbackslash caption\{...\} 定义了图表的标题。如果使用了它，LaTeX 会给你的图表添加「Figure」开头的序号。你可以使用 \textbackslash listoffigures 来生成一个图表的目录。

\textbackslash label\{...\} 创建了一个可以供你引用的标签。

若图或表中有附注，采用英文小写字母顺序编号，附注写在图或表的下方。
国外的期刊习惯将图表的标题和说明文字写成一段，需要改写为标题只含图表的名称，其他说明文字以注释方式写在图表下方，或者写在正文中。

\subsection{多个图片}

如果一个图由两个或两个以上分图组成时，各分图分别以 (a)、(b)、(c)...... 作为图序，并须有分图题。
推荐使用 \pkg{subcaption} 宏包来处理， 比如图~\ref{fig:multi-image}中的图~\ref{fig:subfig-a} 和图~\ref{fig:subfig-b}和。

例如：

\begin{python}
\begin{figure}
	\centering
	\subcaptionbox{分图 A\label{fig:subfig-a}}
	{\includegraphics[width=0.35\linewidth]{example-image-a.pdf}}
	\subcaptionbox{分图 B\label{fig:subfig-b}}
	{\includegraphics[width=0.35\linewidth]{example-image-b.pdf}}
	\caption{多个分图的示例}
	\label{fig:multi-image}
\end{figure}
\end{python}

效果如下：

\begin{figure}
	\centering
	\subcaptionbox{分图 A\label{fig:subfig-a}}
	{\includegraphics[width=0.35\linewidth]{example-image-a.pdf}}
	\subcaptionbox{分图 B\label{fig:subfig-b}}
	{\includegraphics[width=0.35\linewidth]{example-image-b.pdf}}
	\caption{多个分图的示例}
	\label{fig:multi-image}
\end{figure}

\subsection{图片引用}

通过在表格中加入\textbackslash label\{fig\_name\}来定义表格名称，在文本中使用\textbackslash ref\{fig\_name\}来引用。

建议在\textbackslash ref前添加空格~以示美观。

例如：

\begin{python}
此处引用图~\ref{fig:multi-image}，及其图~\ref{fig:subfig-b}和图~\ref{fig:subfig-b}。
\end{python}

效果为：

此处引用图~\ref{fig:multi-image}，及其图~\ref{fig:subfig-b}和图~\ref{fig:subfig-b}。